% ===========================================================================
%  Presentacion
%  Cuentas Nacionales de Transferencia para Mexico, 2016-2024
%
%  Compilar:  pdflatex presentacion   (dos pasadas)
%  Las figuras las genera paper.ipynb; se copian a figures/ desde paper/figures.
% ===========================================================================
\documentclass[aspectratio=169,11pt]{beamer}

\usepackage[utf8]{inputenc}
\usepackage[T1]{fontenc}
\usepackage[spanish,es-nodecimaldot]{babel}
\usepackage{lmodern}
\usepackage{amsmath,amssymb}
\usepackage{booktabs}
\usepackage{graphicx}
\usepackage{array}

\usetheme{default}
\usecolortheme{default}
\setbeamertemplate{navigation symbols}{}
\setbeamertemplate{caption}{\insertcaption}

\definecolor{acento}{HTML}{2C6FBB}
\definecolor{verde}{HTML}{2E8B57}
\definecolor{naranja}{HTML}{E06A1B}
\definecolor{gris}{HTML}{555555}

\setbeamercolor{title}{fg=acento}
\setbeamercolor{frametitle}{fg=acento}
\setbeamercolor{structure}{fg=acento}
\setbeamerfont{frametitle}{size=\large,series=\bfseries}
\setbeamertemplate{itemize item}{\textcolor{acento}{\textbullet}}
\setbeamertemplate{itemize subitem}{\textcolor{gris}{--}}
\setbeamertemplate{footline}{%
  \hfill{\scriptsize\textcolor{gris}{\insertframenumber}}\hspace{1.2em}\vspace{0.6em}}

\graphicspath{{figures/}}

% Subtitulo gris bajo el titulo de lamina
\newcommand{\bajada}[1]{{\small\textcolor{gris}{#1}}\vspace{0.5em}}
% Mensaje de cierre de lamina
\newcommand{\cierre}[1]{%
  \vfill{\small\textcolor{acento}{\textbf{#1}}}}

\title{\textbf{Cuentas Nacionales de Transferencia para México}}
\subtitle{Quién consume, quién aporta y quién financia la diferencia, 2016--2024}
\author{Juan Pablo López Reynosa \and Héctor Villarreal}
\date{\today}

\begin{document}

% ===========================================================================
\begin{frame}[plain]
\titlepage
\end{frame}

% ===========================================================================
\begin{frame}{Objetivo}
\bajada{Qué se construyó y para qué}

\begin{itemize}\itemsep0.7em
  \item Medir el \textbf{déficit del ciclo de vida} en México por edad y por sexo: cuánto consume
        cada edad, cuánto ingreso laboral genera y quién financia la diferencia.
  \item Construir un procedimiento \textbf{reproducible y escalable} que concilie la encuesta de
        hogares con el sistema de cuentas nacionales.
  \item Cubrir cinco levantamientos bienales de la ENIGH, 2016 a 2024, con 1.46 millones de
        registros persona-año.
\end{itemize}

\cierre{La encuesta aporta la forma por edad; las cuentas nacionales aportan el nivel.}
\end{frame}

% ===========================================================================
\begin{frame}{El procedimiento en una vista}
\begin{center}
\includegraphics[width=0.97\linewidth]{fig_pipeline.pdf}
\end{center}
\end{frame}

% ===========================================================================
\section{Resultados}

\begin{frame}{El hallazgo: dos ciclos de vida distintos}
\bajada{El perfil masculino cierra el ciclo; el femenino no lo cierra en ninguna edad}

\begin{center}
\includegraphics[width=0.92\linewidth]{fig_lcd_sexo.pdf}
\end{center}

\vspace{-0.3em}
{\footnotesize El consumo incluye, por convención del marco, el gasto de los hogares \emph{y} el
consumo de gobierno. El déficit del promedio nacional refleja esa definición y la agregación entre
perfiles heterogéneos; no es un juicio sobre una edad en particular.}

\cierre{Las mujeres concentran 76.3\,\% del déficit agregado.}
\end{frame}

% ---------------------------------------------------------------------------
\begin{frame}{Ingreso laboral}
\bajada{La brecha por sexo domina el perfil en todas las edades activas}

\begin{center}
\includegraphics[width=0.92\linewidth]{fig_yl_sexo.pdf}
\end{center}

\cierre{El ingreso laboral femenino pasó de 41.6\,\% a 47.9\,\% del masculino entre 2016 y 2024:
convergencia gradual, sin cambio de orden de magnitud.}
\end{frame}

% ---------------------------------------------------------------------------
\begin{frame}{Consumo y déficit}
\bajada{El déficit es positivo y sostenido; lo que cambia es dónde se concentra}

\begin{center}
\includegraphics[width=0.92\linewidth]{fig_c_lcd.pdf}
\end{center}

\cierre{En términos reales el déficit agregado no crece. El de los mayores de 60 años sube
71.7\,\% por persona, y su peso poblacional pasa de 11.3\,\% a 15.1\,\%.}
\end{frame}

% ---------------------------------------------------------------------------
\begin{frame}{Financiamiento por grupo de edad}
\bajada{Los dos extremos del ciclo se financian por vías opuestas}

\begin{center}
\includegraphics[width=0.92\linewidth]{fig_financiamiento_edad.pdf}
\end{center}

\cierre{Niñez: 88.6\,\% con recursos privados. Vejez: 54.5\,\% con transferencias públicas netas.}
\end{frame}

% ---------------------------------------------------------------------------
\begin{frame}{Jóvenes, mujeres y adultos mayores}
\bajada{La desigualdad por sexo también aparece del lado del financiamiento}

\begin{center}
\includegraphics[width=0.95\linewidth]{fig_financiamiento_grupos.pdf}
\end{center}

\cierre{El déficit de las mujeres de 60 años o más duplica al de los hombres y reciben 24.0\,\%
menos en transferencias públicas netas.}
\end{frame}

% ---------------------------------------------------------------------------
\begin{frame}{El saldo intergeneracional}
\bajada{Pesos constantes de 2024, trimestrales per cápita}

\begin{center}
\begin{tabular}{lrrr}
\toprule
 & 2016 & 2024 & Cambio real \\
\midrule
Transferencia pública neta a 60 años o más & 36\,543 & 37\,943 & +3.8\,\% \\
Pago público neto de la población de 25 a 59 & 9\,317 & 12\,960 & +39.1\,\% \\
\addlinespace
Peso de la población de 60 años o más & 11.3\,\% & 15.1\,\% & $+3.8$ pp \\
\bottomrule
\end{tabular}
\end{center}

\vspace{0.6em}
El pago neto de la edad activa no creció porque le subieran los impuestos ---su pago bruto per
cápita subió 15\,\% real--- sino porque su participación en las transferencias públicas se redujo
a la mitad. El sistema las desplazó hacia las edades mayores, y la base demográfica que lo sostiene
se estrecha.

\cierre{Un cambio demográfico que aumente el peso de la vejez se traduce de forma directa en
presión presupuestaria.}
\end{frame}

% ===========================================================================
\section{Manejo de datos}

\begin{frame}{Fuentes y agregados de control}
\bajada{Toda serie macro utilizada es observada}

\begin{columns}[T]
\begin{column}{0.47\textwidth}
\textcolor{verde}{\textbf{Sistema de Cuentas Nacionales}}
\begin{itemize}\small\itemsep0.35em
  \item Cuentas por sectores institucionales
  \item Remuneración de asalariados e ingreso mixto
  \item Excedente de explotación, consumo de capital fijo y ahorro neto
  \item Consumo final de hogares y de gobierno
  \item Prestaciones sociales, transferencias en especie y del exterior
  \item Impuestos al ingreso y contribuciones sociales
\end{itemize}
\end{column}
\begin{column}{0.47\textwidth}
\textcolor{acento}{\textbf{ENIGH}}
\begin{itemize}\small\itemsep0.35em
  \item Población, ingresos, gasto del hogar y de las personas, trabajos
  \item Una fila por persona y año
  \item 257\,805 a 315\,743 registros por levantamiento
  \item Población expandida de 120.9 a 130.3 millones
\end{itemize}
\end{column}
\end{columns}

\cierre{Las contribuciones sociales imputadas entran al pago público por simetría: el ingreso
laboral las contiene y las prestaciones recibidas incluyen las que financian.}
\end{frame}

% ---------------------------------------------------------------------------
\begin{frame}{La identidad contable y sus anclas}
\bajada{Cada variable micro cierra contra un agregado del sistema de cuentas}

\[
\underbrace{C(a,t) - Y_L(a,t)}_{\text{déficit del ciclo de vida}}
=
\underbrace{\big(\tau^{+}(a,t) - \tau^{-}(a,t)\big)}_{\text{transferencias netas}}
+ \underbrace{\big(Y_K(a,t) - S(a,t)\big)}_{\text{reasignación por activos}}
\]

\vspace{0.3em}
\begin{center}\small
\begin{tabular}{ll}
\toprule
Variable & Agregado de control \\
\midrule
Ingreso laboral & Remuneración $+\ \tfrac{2}{3}$ del ingreso mixto neto \\
Reasignación por activos & Ingreso de activos ampliado, neto de ahorro \\
Consumo total, educación y salud & Consumo final de hogares y gobierno \\
Transferencias públicas recibidas & Prestaciones sociales y transferencias en especie \\
Transferencias del exterior & Transferencias netas del exterior \\
\bottomrule
\end{tabular}
\end{center}

\vspace{0.4em}
{\footnotesize La división del ingreso mixto en dos tercios laborales y un tercio de capital es la
convención del manual de Naciones Unidas, aplicada por igual en el micro y en el macro.}
\end{frame}

% ---------------------------------------------------------------------------
\begin{frame}{Cómo se concilian micro y macro}
\bajada{Un único factor por variable y año}

\[
\lambda_v(t) = \frac{V_{macro}(t)}{\sum_i x_{iv}(t)\cdot f_i \cdot 4}
\qquad\Longrightarrow\qquad
x^{adj}_{iv}(t) = x_{iv}(t)\cdot\lambda_v(t)
\]

\begin{itemize}\itemsep0.5em
  \item \textbf{Cierre exacto por construcción}: ninguna imputación previa puede mover un agregado
        anual.
  \item \textbf{La forma no cambia}: se preserva toda la información distributiva y etaria de la
        encuesta; sólo cambia la escala vertical.
  \item \textbf{El factor de expansión no se toca}: la conciliación no altera la estructura
        demográfica.
\end{itemize}

\cierre{Los factores van de 1.9 a 9.9: la encuesta capta entre 10\,\% y 53\,\% del agregado
correspondiente, según la variable.}
\end{frame}

% ---------------------------------------------------------------------------
\begin{frame}{Asignación del gasto del hogar}
\bajada{La encuesta observa gasto del hogar; el perfil por edad exige repartirlo entre personas}

\begin{itemize}\itemsep0.5em
  \item Algoritmo iterativo de \textbf{punto fijo} que construye su propia escala de equivalencia
        por edad a partir del dato, en vez de importar una escala externa.
  \item Parte de un reparto per cápita, estima la participación de cada edad en el agregado,
        redistribuye según ese perfil y repite hasta convergencia.
  \item Conserva el total del hogar de forma exacta en cada iteración: ninguna configuración del
        algoritmo puede mover un agregado.
\end{itemize}

\cierre{La alternativa habitual ---imponer una escala construida para otro país--- se evita por
diseño.}
\end{frame}

% ---------------------------------------------------------------------------
\begin{frame}{Imputación fiscal}
\bajada{La encuesta reporta ingreso neto, no carga fiscal}

\begin{itemize}\itemsep0.5em
  \item \textbf{Contribuciones sociales}: salario base de cotización con piso y topes por régimen.
        Entra la contribución patronal y la de gobierno; la del trabajador ya viene descontada del
        ingreso observado.
  \item \textbf{Impuesto sobre la renta}: la tarifa se aplica al ingreso bruto y la encuesta
        observa el neto, de modo que el bruto se recupera por iteración de punto fijo.
  \item Se añade el tramo de \textbf{personas morales} sobre el ingreso de capital observado, sin
        el alquiler imputado por no ser hecho gravable.
  \item Todo se calcula sobre \textbf{base anual}: el salario base, sus topes y la tarifa
        progresiva son parámetros anuales.
\end{itemize}

\cierre{La construcción se apoya en el Simulador Fiscal del CIEP, tanto en las tarifas como en el
criterio de base gravable.}
\end{frame}

% ---------------------------------------------------------------------------
\begin{frame}{Transferencias dentro del hogar}
\bajada{Hacer visible una reasignación que suele quedar implícita en un residuo}

\begin{itemize}\itemsep0.5em
  \item Para cada persona se calcula su ingreso propio y su posición de déficit o superávit.
  \item Dentro de cada hogar se emparejan miembros deficitarios con superavitarios, y el saldo
        neto \textbf{suma exactamente cero por hogar}.
  \item A diferencia de un residuo contable, el método \textbf{puede dejar déficit sin cubrir}
        cuando el hogar no tiene superávit suficiente: ese remanente se financia por fuerza desde
        fuera del hogar.
  \item El bloque privado queda observado, sin reescalar, porque el sistema de cuentas no publica
        un agregado de transferencias privadas netas.
\end{itemize}

\cierre{El gobierno aparece como mediador entre generaciones, no como agente: recauda de unas
edades y transfiere a otras.}
\end{frame}

% ===========================================================================
\section{Cierre}

\begin{frame}{Limitaciones declaradas}
\bajada{Nueve limitaciones documentadas, con su magnitud cuantificada}

\begin{itemize}\small\itemsep0.45em
  \item \textbf{Impuesto sobre la renta.} El agregado de control es la recaudación total del
        impuesto al ingreso. Lo que la encuesta permite imputar ---el impuesto de quienes reciben
        sueldo y el de las rentas de capital que declaran los hogares--- equivale apenas a un
        tercio de esa recaudación. La diferencia corresponde sobre todo a utilidades de empresas
        que nunca llegan al hogar como ingreso observado, de modo que el nivel del impuesto por
        edad queda sobreestimado aunque su forma sea informativa.
  \item \textbf{Remesas.} La encuesta registra apenas una novena parte del monto que reportan las
        cuentas nacionales, y lo hace sobre muy pocos hogares receptores. Al ajustar el total, todo
        el faltante se carga sobre el monto que recibe cada hogar y nada sobre el número de hogares
        que reciben, cuando lo probable es que ambos estén subregistrados. La distribución por edad
        y sexo del canal es confiable; su nivel no.
  \item El consumo de gobierno se imputa plano: no hay fuente de utilización de servicios públicos
        por edad, y el gasto privado en salud puede ser inverso al beneficio público recibido.
  \item El reescalamiento proporcional no corrige subcaptación diferencial por edad o por decil.
\end{itemize}

\cierre{El procedimiento sirve como control macroeconómico de los perfiles por edad, no como
instrumento para medir desigualdad del ingreso de capital.}
\end{frame}

% ---------------------------------------------------------------------------
\begin{frame}{Reproducibilidad y escalabilidad}
\bajada{Qué hace utilizable al procedimiento más allá de este ejercicio}

\begin{itemize}\itemsep0.5em
  \item \textbf{Cada decisión es explícita y sustituible}: mejorar la identificación de un bloque
        no obliga a rehacer los demás.
  \item \textbf{Pruebas de cierre reproducibles}: que ninguna imputación mueva un agregado anual,
        y que los perfiles por edad reproduzcan exactamente el resultado de la base persona.
  \item \textbf{La brecha se reporta, no se fuerza}: el déficit calculado por sus dos lados no
        coincide de forma exacta, y la diferencia es analíticamente igual al bloque privado
        observado.
  \item \textbf{Todo se entrega dentro de un repositorio}: datos de control, cuadernos, parámetros
        fiscales y documentos, con la cadena que los produce y su documentación de operación.
\end{itemize}
\end{frame}

% ---------------------------------------------------------------------------
\begin{frame}{Siguiente paso}
\bajada{Qué habilita la base construida}

\begin{itemize}\itemsep0.7em
  \item El producto es una base con observaciones a \textbf{nivel de agente representativo},
        construida desde la encuesta y conciliada con las cuentas nacionales.
  \item Esa estructura permite \textbf{proyectar los perfiles} hacia adelante bajo distintos
        escenarios demográficos y económicos.
  \item Y evaluar el efecto de cambios de política sobre cada grupo de edad y sexo, en vez de
        sobre el promedio nacional.
\end{itemize}

\vfill
\begin{center}
{\large\textcolor{acento}{\textbf{Gracias}}}
\end{center}
\end{frame}

% ===========================================================================
\appendix

\begin{frame}{Respaldo: alquiler imputado de la vivienda propia}
\begin{itemize}\small\itemsep0.45em
  \item La encuesta capta gasto como desembolso: un hogar propietario consume servicio
        habitacional sin desembolsar, y ese servicio no aparece.
  \item El marco lo reconoce como dos cosas a la vez: consumo del servicio e ingreso del activo.
        Entra por ambos lados de la identidad.
  \item El reparto dentro del hogar es per cápita por los dos lados, siguiendo el criterio del
        Simulador Fiscal del CIEP.
  \item Tiene contraparte macro ---el excedente de operación de los hogares--- pero no agregado de
        control propio.
\end{itemize}
\end{frame}

\begin{frame}{Respaldo: transferencias del exterior}
\begin{itemize}\small\itemsep0.45em
  \item La encuesta capta entre 10.1\,\% y 12.0\,\% del agregado de control, sobre una base de
        1.6\,\% de las personas.
  \item El reescalamiento lleva el monto medio por receptor de 16\,088 a 145\,339 pesos
        trimestrales.
  \item La \textbf{forma} del canal por edad y por sexo proviene de la encuesta y es robusta al
        ajuste; el \textbf{nivel} depende del agregado de control.
  \item Corregirlo exige repartir el ajuste entre número de receptores y monto por receptor, con
        un objetivo externo de incidencia.
\end{itemize}
\end{frame}

\begin{frame}{Respaldo: composición de la muestra}
\begin{center}\small
\begin{tabular}{lrrrrr}
\toprule
Año & Registros & Hogares & Población & Edad media & 60 años o más \\
\midrule
2016 & 257\,805 & 70\,311 & 120.9 & 31.0 & 11.3\,\% \\
2018 & 269\,206 & 74\,647 & 123.9 & 31.8 & 12.1\,\% \\
2020 & 315\,743 & 89\,006 & 126.8 & 33.2 & 13.6\,\% \\
2022 & 309\,684 & 90\,102 & 129.0 & 33.6 & 14.2\,\% \\
2024 & 308\,598 & 91\,414 & 130.3 & 34.4 & 15.1\,\% \\
\bottomrule
\end{tabular}
\end{center}
\vspace{0.5em}
{\footnotesize Población en millones de personas expandidas. Fuente: elaboración propia con ENIGH
(INEGI).}
\end{frame}

\end{document}
